% ══════════════════════════════════════════════════════════════════════════════
%  2027 WORKBOOK STYLE (Year 7 workbook redesign, Nov 2026) — A4 version
%  The look of the A5 trial (space/booklet/a5-test) on an A4 page:
%    · colour cover from the unit's cover art, black-and-white inside
%    · headings in Outfit (the site's heading face), body in Lato
%    · each experiment opens with a band: number, title, ink symbol, heavy rule
%    · boxes told apart by frame and icon, never by colour:
%        I do ●  solid · We do ●● dashed · You do ●●● heavy
%        Key idea / technique ◆ thin · Above level ▲ double · Country ◎ dotted
%    · skill tags are ink tabs carrying the code (QP1, PC4 …) plus the name
%  Everything above this block still names the old colours; they are mapped to
%  ink here, so the body never needs to know which look it is printed in.
% ══════════════════════════════════════════════════════════════════════════════

\usetikzlibrary{patterns}
% ── Fonts ───────────────────────────────────────────────────────────────────
% Bold is Outfit SemiBold, the weight the site uses; Outfit has no italic,
% so a slanted one stands in.
\renewfontfamily\headingfont{Outfit}[
  Path=./fonts/, Extension=.ttf, Ligatures=TeX,
  UprightFont=Outfit-Regular, BoldFont=Outfit-SemiBold,
  ItalicFont=Outfit-Regular, ItalicFeatures={FakeSlant=0.18},
  BoldItalicFont=Outfit-SemiBold, BoldItalicFeatures={FakeSlant=0.18} ]

% ── Black and white inside ──────────────────────────────────────────────────
\definecolor{Ink}{HTML}{111111}
\colorlet{NavyBlue}{Ink}\colorlet{MidBlue}{Ink}\colorlet{TealGreen}{Ink}
\colorlet{GoldDark}{Ink}\colorlet{PurpleDark}{Ink}\colorlet{OrangeDark}{Ink}
\colorlet{TableHead}{Ink}
\colorlet{Blue}{Ink}\colorlet{Teal}{Ink}\colorlet{Amber}{Ink}\colorlet{Coral}{Ink}
\colorlet{Violet}{Ink}\colorlet{Sky}{Ink}
\colorlet{LightBlue}{white}\colorlet{LightGreen}{white}\colorlet{GoldLight}{white}
\colorlet{PurpleLight}{white}\colorlet{OrangeLight}{white}\colorlet{RowAlt}{white}
\colorlet{LightGrey}{white}\colorlet{Paper}{white}
\colorlet{DarkGrey}{black!72}\colorlet{MidGrey}{black!45}\colorlet{Line}{black!45}
\arrayrulecolor{black!60}

% ── Page style ──────────────────────────────────────────────────────────────
\fancypagestyle{fancy}{%
  \fancyhf{}%
  \renewcommand{\headrulewidth}{0.4pt}%
  \renewcommand{\headrule}{\hbox to\headwidth{\color{Ink}\leaders\hrule height \headrulewidth\hfill}}%
  \fancyhead[L]{\headingfont\footnotesize \unitname\ \textbullet\ Experiment Booklet}%
  \fancyhead[R]{\headingfont\footnotesize\leftmark}%
  \fancyfoot[C]{\headingfont\footnotesize\thepage}}
\fancypagestyle{plain}{\fancyhf{}\renewcommand{\headrulewidth}{0pt}%
  \fancyfoot[C]{\headingfont\footnotesize\thepage}}
\fancypagestyle{opener}{\fancyhf{}\renewcommand{\headrulewidth}{0pt}%
  \fancyfoot[C]{\headingfont\footnotesize\thepage}}
\pagestyle{fancy}
\markboth{}{}

\titleformat{\subsubsection}{\headingfont\normalsize\bfseries}{}{0em}{}

% ── Box system: told apart by frame and icon, never by colour ───────────────
\newcommand{\icodots}[1]{\tikz[baseline=-0.6ex]{\foreach \k in {1,...,#1}{\fill[white] (\k*1.9mm,0) circle (0.7mm);}}}
\newcommand{\icokey}{\tikz[baseline=-0.6ex]{\fill[white] (0,0.9mm)--(0.9mm,0)--(0,-0.9mm)--(-0.9mm,0)--cycle;}}
\newcommand{\icoup}{\tikz[baseline=-0.6ex]{\fill[white] (-1mm,-0.8mm)--(1mm,-0.8mm)--(0,1mm)--cycle;}}
\newcommand{\icocountry}{\tikz[baseline=-0.6ex]{\draw[white,line width=0.5pt] (0,0) circle (1mm); \fill[white] (0,0) circle (0.45mm);}}
\newcommand{\icoskill}{\tikz[baseline=-0.6ex]{\fill[white] (-0.85mm,-0.85mm) rectangle (0.85mm,0.85mm);}}
\newcommand{\tagicon}[1]{\tikz[baseline=-0.6ex]\node[fill=Ink, rounded corners=0.8mm, inner xsep=1.2mm, inner ysep=1mm]{#1};}
\tcbset{wbbase/.style={enhanced, colback=white, arc=1.2mm,
    left=3.5mm, right=3.5mm, top=4.2mm, bottom=2.6mm, before skip=11pt, after skip=8pt,
    attach boxed title to top left={xshift=3mm, yshift=-2.4mm},
    boxed title style={colback=Ink, colframe=Ink, arc=0.8mm, boxrule=0pt,
                       left=1.6mm, right=1.8mm, top=0.6mm, bottom=0.6mm},
    fonttitle=\headingfont\bfseries\footnotesize, coltitle=white}}

% I do: the worked example (solid frame, one dot)
\renewtcolorbox{workedexbox}[1]{wbbase, breakable, colframe=Ink, boxrule=1pt,
  title={\icodots{1}\ I do \textbullet\ #1}}
\renewenvironment{workedex}[1]
  {\refstepcounter{wenum}\label{we:\thewenum}\begin{workedexbox}{Worked example \thewenum\ --- #1}}
  {\end{workedexbox}}
% We do: straight after its worked example (dashed frame, two dots)
\renewtcolorbox{wedobox}[1]{wbbase, breakable=false, frame hidden,
  borderline={1.1pt}{0pt}{Ink, dash pattern=on 4pt off 2.6pt},
  title={\icodots{2}\ #1}}
\renewenvironment{wedo}[1]
  {\par\begin{wedobox}{We do \thewenum\ --- #1}\label{wd:\thewenum}\ignorespaces}
  {\end{wedobox}}
% You do: a set of practice tasks on your own (heavy frame, three dots)
\newtcolorbox{youdobox}[1][Check your understanding]{wbbase, breakable=false,
  colframe=Ink, boxrule=2.4pt, top=5mm,
  title={\icodots{3}\ You do \textbullet\ #1}}
% Key idea, technique, safety, word banks (thin frame, diamond)
\renewtcolorbox{infobox}[1]{wbbase, breakable, colframe=Ink, boxrule=0.5pt,
  title={\icokey\ #1}}
\renewtcolorbox{scaffold}[1]{wbbase, breakable, colframe=Ink, boxrule=0.5pt,
  title={\icokey\ #1}}
\renewtcolorbox{technique}[2]{wbbase, breakable, colframe=Ink, boxrule=0.5pt,
  title={\icokey\ Technique \textbullet\ #1\enspace\textbar\enspace #2 \skname{#2}}}
% Above level (double frame, triangle)
\renewtcolorbox{yr8box}{wbbase, breakable, frame hidden,
  borderline={0.6pt}{0pt}{Ink}, borderline={0.6pt}{1.6pt}{Ink},
  title={\icoup\ Going further \textbullet\ Above level}}
% Country (dotted frame, ring)
\newtcolorbox{countrybox}[1]{wbbase, breakable, frame hidden,
  borderline={1.4pt}{0pt}{Ink, line cap=round, dash pattern=on 0pt off 2.4pt},
  title={\icocountry\ #1}}
% Learning intentions and skills focus (thin frame; square for skills)
\renewtcolorbox{lipanel}{wbbase, breakable=false, colframe=Ink, boxrule=0.5pt,
  title={\icokey\ Learning intentions}}
\renewtcolorbox{skillsfocus}{wbbase, breakable=false, colframe=Ink, boxrule=1pt,
  title={\icoskill\ Skills focus \textbullet\ the science skills this experiment builds},
  before upper={\begin{multicols}{2}\raggedcolumns}, after upper={\end{multicols}}}

% ── Skill tags: an ink tab with the code, then the name ─────────────────────
\makeatletter
\renewcommand{\skc}[1]{%
  \tikz[baseline=(c.base)]\node[rounded corners=0.8mm, fill=Ink, text=white, inner xsep=1.5mm, inner ysep=0.7mm,
    font=\headingfont\bfseries\scriptsize] (c) {#1};}
\renewcommand{\sk}[1]{\mbox{\skc{#1}\,{\headingfont\bfseries\footnotesize\skname{#1}}}}
\makeatother
\renewcommand{\focus}[2]{\par\noindent\makebox[1.1cm][l]{\skc{#1}}\parbox[t]{\dimexpr\linewidth-1.1cm\relax}{\footnotesize\RaggedRight\textbf{\skname{#1}}\,---\,#2}\par\vspace{3pt}}

% ── Experiment opener band ──────────────────────────────────────────────────
% \experimentdiv{title}{summary}{lesson}: number, title and the experiment's ink
% symbol (\expsymbol{N}), then a heavy rule. Steps the section counter, so
% tasks, worked examples and figures are N.M.
\newcommand{\expword}{Experiment}\newcommand{\expcaps}{EXPERIMENT}\newcommand{\symid}{\thesection}
\renewcommand{\experimentdiv}[3]{%
  \clearpage\thispagestyle{opener}%
  \refstepcounter{section}\markboth{\expword\ \thesection\ \textbullet\ #1}{}%
  \pdfbookmark[1]{\expword\ \thesection{} — #1}{exp.\thesection}%
  \noindent\begin{minipage}[b]{\dimexpr\linewidth-46mm\relax}
    {\headingfont\bfseries\footnotesize\addfontfeatures{LetterSpace=14}\unitcaps\ \textbullet\ \expcaps}\par\vspace{1mm}
    {\headingfont\bfseries\fontsize{40}{40}\selectfont\thesection}\par\vspace{1mm}
    {\headingfont\bfseries\fontsize{19}{22}\selectfont\RaggedRight #1\par}\vspace{1mm}
    {\small\itshape #2}\par\vspace{1mm}
    {\headingfont\footnotesize Lesson #3}
  \end{minipage}\hfill
  \begin{minipage}[b]{44mm}\centering\resizebox{!}{27mm}{\expsymbol{\symid}}\end{minipage}\par
  \vspace{2mm}\noindent\rule{\linewidth}{1.6pt}\par\vspace{1mm}}
\renewcommand{\lessoncodes}[1]{{\small\itshape VC 2.0: #1}\par\vspace{2pt}}

% Part heading inside an experiment: every part starts a new page; the
% starred form follows straight on.
\RenewDocumentCommand{\exppart}{s m m}{%
  \par\IfBooleanTF{#1}{\needspace{10\baselineskip}\vspace{8pt}}{\clearpage}%
  {\headingfont\bfseries\LARGE #2}\par\vspace{1mm}\noindent\rule{\linewidth}{0.6pt}\par\vspace{2mm}}

% Key vocabulary: a heading like the trial's, then the write-in table
\renewcommand{\vocabhead}{%
  \par\vspace{4pt}{\headingfont\bfseries\Large Key vocabulary}\par\vspace{1pt}%
  \noindent{\small Copy the meaning of each word from the slide. The first one has been done for you.}\par}

% Table headings: bold Outfit on white, the grid in ink
\renewcommand{\thead}[1]{\headingfont\bfseries #1}
\renewcommand{\tickbox}{\tikz[baseline=-0.6ex]\draw[Ink, line width=0.7pt, rounded corners=1pt] (0,-0.17) rectangle (0.34,0.17);}

% Glossary rows: \glossrow{word}{label}
\newcommand{\glossrow}[2]{\rule{0pt}{5mm}\textbf{#1} & \pageref{#2} \\ \hline}

% ── Ink symbols for the experiment openers ───────────────────────────────────
% One diagrammatic idea per experiment, drawn in ink from the same 2D
% equipment shapes the booklet teaches.
\tikzset{sym/.style={line width=1.5pt, draw=Ink, line cap=round, line join=round}}
% The openers use the unit's own lesson symbols (art/lN-ink.pdf, from
% tools/forces_symbols.py): Experiments 1-5 = Lessons 3, 6, 8, 10, 13.
\newcommand{\expsymbol}[1]{\edef\exparg{\ifcase#1\or3\or6\or8\or10\or13\fi}\includegraphics[height=27mm]{art/l\exparg-ink.pdf}}
